\chapter*{Introduction générale}
\addcontentsline{toc}{chapter}{Introduction générale}
\markboth{Introduction générale}{}

Les entreprises produisent aujourd'hui de très grandes quantités de données issues de leurs systèmes
opérationnels : ventes, achats, stocks, paiements, logistique. Ces données constituent une ressource
stratégique pour suivre l'activité, mesurer la performance et éclairer la prise de décision. Leur
exploitation demeure pourtant difficile : elles sont dispersées dans des bases transactionnelles conçues
pour l'enregistrement des opérations et non pour l'analyse, et leur interrogation exige une maîtrise
technique, en particulier du langage SQL, que la plupart des décideurs ne possèdent pas.

Deux évolutions récentes changent profondément cette situation. D'une part, l'émergence de la
\emph{Modern Data Stack} --- entrepôts de données dans le cloud, outils d'ingestion déclaratifs,
transformations versionnées et testées, orchestrateurs orientés données --- permet d'industrialiser la
chaîne de valeur de la donnée avec les pratiques du génie logiciel. D'autre part, les progrès des
grands modèles de langage (LLM) rendent possible l'interrogation des données en langage naturel : une
question posée en français peut être traduite automatiquement en une requête SQL exécutable.

C'est dans ce double contexte que s'inscrit ce projet, réalisé dans le cadre du Master II Informatique,
option Business Intelligence, de l'Université Cheikh Anta Diop de Dakar. Il vise à concevoir et à mettre
en œuvre une \textbf{Modern Data Platform} reposant sur une \textbf{architecture Médaillon} (Bronze,
Silver, Gold) dans Snowflake, alimentée depuis SQL Server par \textbf{DLT}, transformée avec
\textbf{dbt} et orchestrée par \textbf{Dagster}, puis à développer une application \textbf{Streamlit}
de type \emph{Text-to-SQL} s'appuyant sur \textbf{Snowflake Cortex AI} pour permettre aux décideurs
d'obtenir leurs indicateurs sans écrire une ligne de SQL.

Le présent rapport est organisé comme suit :
\begin{itemize}
  \item le \textbf{chapitre~\ref{ch:contexte}} présente le contexte, la problématique, les objectifs et le
        jeu de données Wide World Importers ;
  \item le \textbf{chapitre~\ref{ch:etatart}} expose les concepts et technologies mobilisés ;
  \item le \textbf{chapitre~\ref{ch:conception}} décrit l'architecture retenue et la conception du modèle
        de données ;
  \item le \textbf{chapitre~\ref{ch:realisation}} détaille la réalisation de la plateforme : ingestion,
        couches Silver et Gold, historisation, qualité des données et orchestration ;
  \item le \textbf{chapitre~\ref{ch:application}} est consacré à l'application Text-to-SQL ;
  \item le \textbf{chapitre~\ref{ch:industrialisation}} traite de l'industrialisation : traitements
        incrémentaux, alertes, déploiement dans Snowflake et intégration continue ;
  \item le \textbf{chapitre~\ref{ch:resultats}} présente les résultats obtenus, les difficultés
        rencontrées et leurs solutions ;
  \item enfin, la conclusion dresse le bilan du projet et en propose les perspectives.
\end{itemize}
